% Changed: argument and layout integration; measured claims retained.
% Protected: existing citations, equations and labels unless explicitly documented.
% Meaning-Check: PRESERVED
% Risk-Flags: none
\section{Implementation and Measurement}
\label{sec:implementation}

\subsection{Integration with the Serving Engine}
The prototype targets the vLLM 0.10.2 V1 scheduler and cache manager. It comprises a request-side compiler and qualification controller, an in-engine verification executor, and PSTP hooks at cache-index publication. Model weights and attention kernels remain unchanged. Candidate identities are derived after full prompt rendering and tokenization. The archived performance clients use policy/template cache identities for verification and compiled requests. A separately evaluated authorization integration extends the identity to every engine route: a trusted authority validates membership and the active epoch, then derives a domain/epoch-bound cache salt for raw, fallback, verifier and compiled requests. Salt propagation through the native block-hash chain makes physical sharing follow the authorized group (Section~\ref{sec:authorization-results}).

The deployment client checks HMAC binding/lifecycle; the query client checks frozen identities and scope (Table~\ref{tab:qualification-protocols}). Both invoke in-engine verification/PSTP. The all-route integration permits authorized raw fallback for an invalid compiled binding, but rejects expired identities before dispatch. It is measured in its own regression cohort; performance tables retain their archived clients.

Compiler and publication overheads measured in the deployment implementation are 0.389 ms mean for CALC, 0.249 ms for the residency lookup, 0.109 ms for the template-scope check, and 0.636 ms for PSTP finalization. The query-qualified studies measure their complete controller path afresh rather than importing these component timings as constants.

\subsection{Task-Preserving Input Construction}
The adapter retains each source in full when it fits the common 16K window. Otherwise a fixed head/tail rule crops optional paragraph/code-line substrings, protecting the complete question and code prefix, original order/whitespace and an omission ledger. Selection never inspects answers, outputs, scores, margins or qualification. Adjacent pieces of one source segment remain contiguous in both routes.

Both routes receive identical retained content, checked against every model/tenant/turn/layout condition, chat template and output reserve. Offline tokenization is compared with API prompt IDs. Common immutable packing is outside the clock; request rendering, planning and verification are included.

The input-construction experiment compares fixed-word-cap, whitespace-normalized packing with structure-preserving rendering under the same window, admission rules and scorer. This isolates how task-critical content affects absolute quality and admissible reuse (Section~\ref{sec:fidelity}).

\subsection{Execution Environment and Timing}
The primary hardware is two RTX 3090 GPUs with 24 GB each. Single-model cells use one GPU; Qwen3-14B uses tensor parallelism across two GPUs in the deployment evaluation. The query-qualified environment is Python 3.10.20, PyTorch 2.8.0+cu128, vLLM 0.10.2, transformers 4.55.2, Triton 3.4.0, and FP16. Model checkpoints, tokenizers, runtime files, inputs, and policy identities are recorded per cell. The primary task-preserving cells use two concurrent requests; ten registered tenants per source describe sharing fan-out, not ten simultaneous GPU requests.

A fresh server process is used for each cell. The closed-loop client runs one tenant pair through ten rounds, then advances to the next pair. The next round becomes ready when the preceding business wave finishes; synchronous qualification performed between waves contributes to subsequent readiness delay. For request $r$,
\begin{equation}
 T^{\rm ready}_{\rm TTFT}(r)=t_{\rm first}(r)-t_{\rm ready}(r).
\end{equation}
Full service wall includes all business responses, first-use raw witnesses, planning, information and verification calls, unsuccessful attempts, waiting, scoring, logging, and final drain. Server startup, common pre-measurement infrastructure checks, configuration loading, and immutable input preparation have separately declared boundaries.

Closed-loop throughput is completed business requests divided by full service wall. The paired multiplier is the raw wall divided by the ContractKV wall for an equal request population. In the separate fixed-arrival deployment studies, SLO goodput counts responses meeting a specified readiness deadline per unit service time. Their arrival traces and workload protocol are reported separately.

Broad deployment measurements use an explicit residual-verification accounting rule,
\begin{equation}
 C_{\rm exposed}=\max(C_{\rm verify}-\Delta_{\rm slack},0),
\label{eq:slack}
\end{equation}
with 2,000 ms slack for Qwen/Mistral and 6,000 ms for the broad Llama protocol. All new query-qualified results use the directly measured complete service clock with zero slack. Component RPC sums provide a work breakdown alongside the directly measured wall clock.

\subsection{Physical Accounting and Reproducibility}
At initial scheduler allocation, actual prompt IDs and native hit count $h_r$ give $n_r-h_r$ newly required prefill tokens, with separate verification counters. These measure work; latency is measured directly.

Source groups and policies are frozen before opening each independent cohort. Source-cluster statistics keep tenant/query repetitions nested and retain process variation. The artifact supplies request-level data, model/runtime manifests, source snapshots, aggregation code and archive checksums.\footnote{\url{https://github.com/zhouz4884-coder/contractkv}. Performance evidence: \texttt{84dd0366}. Protocol-specific freezes, additional evidence and hashes are in the supplement register.}
